\documentclass[12pt]{article}
\usepackage[utf8]{inputenc}
\usepackage[T1]{fontenc}
\usepackage{hyperref}
\title{Arabic, Judeo-Iraqi Language Resources for K-12 Curriculum Development}
\author{Compiled by Claude Sonnet 5.5 (claude-sonnet-5-5) for Andr\'as Kornai}
\date{}
\begin{document}
\maketitle

\section{Introduction}
Judeo-Iraqi Arabic (ISO 639-3 yhd; Glottocode jude1266) is an endangered Afro-Asiatic language. It is used as a first language by older adults only and is not known to be taught in schools. The web search found very little, namely one online resource. Nothing else (grammars, textbooks, school materials) turned up, so this report is short.

\section{Primary Educational Resources}
No primers, textbooks or school programmes were found in the search results.

\section{Linguistic and Reference Materials}
\begin{itemize}
\item Judeo-Iraqi Arabic Living Dictionary (Living Dictionaries). It offers Hebrew and Arabic orthography options. Entries give English and Hebrew translations, part of speech and semantic domain. It is part of Heirloom, a language reclamation programme run by the Jewish Language Project and the Oxford School of Rare Jewish Languages. \url{https://livingdictionaries.app/judeo-iraqi-arabic}
\end{itemize}
The site menu also shows a ``Grammar'' section. The search results did not show its content, so its coverage is unconfirmed.

\section{Religious and Community Resources}
No religious texts or community organisations were found in the search results, beyond the Heirloom programme named above.

\section{Digital and Audio Resources}
The Living Dictionary above is the only digital resource found. Audio content was not confirmed.

\section{Curriculum Development Context}
With only older adults as speakers and almost no documented material, any curriculum would be a revitalization or heritage effort. It would have to begin by documenting speakers and expanding the Living Dictionary. Further searching for scholarly grammars and archives is needed.

\end{document}